\documentclass[10pt,a4paper]{article}
\usepackage[margin=1in]{geometry}
\usepackage{booktabs}
\usepackage{xcolor}
\usepackage{listings}
\usepackage{hyperref}
\usepackage{microtype}

\lstset{
  basicstyle=\ttfamily\small,
  backgroundcolor=\color{black!4},
  frame=single, framesep=4pt, framerule=0pt,
  breaklines=true, columns=flexible,
}

\setlength{\parindent}{0pt}
\setlength{\parskip}{0.5em}

\title{\vspace{-2em}\textbf{Experiment G --- Quadrature Speed/Equivariance Pareto}}
\author{}
\date{}

\begin{document}
\maketitle
\vspace{-3em}

What is the right way to compare DH vs GL? \textbf{Same equivariance error},
then ask how much wall-clock time each method needs to reach it.

\section*{Setup (fairchem default EquiformerV2, lmax$=$6, mmax$=$2)}

Real \texttt{EquiformerV2Backbone} from fairchem: 12 layers, 128 channels,
24 \texttt{S2Activation} modules. NVIDIA A100-SXM4-40GB, idle.

These are the function-signature defaults of \texttt{EquiformerV2}
in fairchem (\texttt{num\_layers=12}, \texttt{sphere\_channels=128},
\texttt{lmax\_list=[6]}, \texttt{mmax\_list=[2]}); the published OC20
SOTA model is larger (20 layers, lmax=8, mmax=8, $\sim$150M params),
so the $+31\%$ savings here is a lower bound for the production model.
The 24 modules $=$ $2 \times$ num\_layers (one S2 act in graph attention
$+$ one in FFN per \texttt{TransBlockV2}).

\section*{Two measurements at the same (lmax, mmax)}

\textbf{Equivariance error} of one S2-Act kernel (random coefficients,
5 rotations $\times$ 10 inputs); SiLU activation. \emph{Source:}
\texttt{verify\_gl\_grid.py}.

\textbf{Wall-clock} for one full forward (QM9, batch$=$8). Protocol:
5 independent runs per config; each run uses 10 warmup $+$ 30 measured
forwards on \emph{distinct} QM9 batches drawn without replacement
(no batch reuse within a run, fresh RNG seed per run).
\texttt{cuda.sync()} around every measurement; fresh model construction
per (run, config). 95\% CI is reported as
$t_{n_\text{runs}-1, 0.975} \times \mathrm{s.e.}(\text{run\_means})$
with $n_\text{runs}=5$ as the unit of replication. \emph{Source:}
\texttt{verify\_speedup\_v2.json}.

\section*{Combined results}

\begin{center}
\small
\begin{tabular}{lcrl}
\toprule
Config & Pts & Equiv err & Fwd (ms, 95\% CI) \\
\midrule
DH default                   & 70  & $4.43\!\times\!10^{-1}$ & $113.50 \pm 1.87$ \\
\textbf{DH 2$\times$}        & 784 & $\mathbf{3.27\!\times\!10^{-1}}$ & $\mathbf{164.66 \pm 1.62}$ ($+45.1\%$, $p<10^{-4}$) \\
GL min                       & 91  & $4.29\!\times\!10^{-1}$ & $112.68 \pm 1.87$ ($-0.7\%$, n.s.) \\
\textbf{GL match-DH}         & 182 & $\mathbf{3.28\!\times\!10^{-1}}$ & $\mathbf{113.17 \pm 2.25}$ ($-0.3\%$, n.s.) \\
\textbf{GL 2$\times$}        & 392 & $\mathbf{3.28\!\times\!10^{-1}}$ & $\mathbf{112.61 \pm 2.35}$ ($-0.8\%$, n.s.) \\
\bottomrule
\end{tabular}
\end{center}

\section*{Matched-equivariance comparison}

For target equiv err $= 3.28 \times 10^{-1}$ (the saturation floor at this mmax cropping):

\begin{itemize}\setlength\itemsep{0pt}
\item Cheapest DH: \textbf{DH 2$\times$ at 164.66 ms}
  (DH default at 70 pts only reaches $4.43\!\times\!10^{-1}$, not enough)
\item Cheapest GL: \textbf{GL match-DH at 113.17 ms}
\item \textbf{Savings: $51.5 \pm 2.36$ ms = $-31.3\%$ time, $p < 10^{-4}$ from
  Welch's t-test on $n=5$ run means}
\end{itemize}

\textbf{So: at matched equivariance accuracy, GL saves $\sim$31\% wall-clock per
forward.}

Equivalently: GL achieves DH 2$\times$'s equivariance at the same wall-clock
as DH default (the GL-vs-DH-default differences are statistically zero).
The $+51.5$ ms premium that DH pays to reach the equivariance floor goes
away under GL.

\section*{Caveats}

\begin{enumerate}\setlength\itemsep{0pt}
\item Result is for this specific (lmax$=$6, mmax$=$2) config and this GPU.
  Different model depths or hardware will scale it.
\item Equivariance is measured on \emph{one} S2-Act kernel with random
  inputs, not on a trained model's intermediate features. Trained-model
  features are smoother (absolute equiv numbers will be much lower); the
  \emph{ranking} between quadratures is what is claimed to transfer.
  See \texttt{dropin\_pretrained.json} for validation on actual
  4-layer trained checkpoints.
\item The kernel-level equivariance gap between GL match-DH and DH 2$\times$
  is within rounding ($3.28\!\times\!10^{-1}$ vs $3.27\!\times\!10^{-1}$).
\item The mmax$=$2 cropping caps achievable equivariance at $\sim$0.328
  regardless of grid; lower error needs mmax$=$lmax (architectural change).
  At lmax$=$mmax both methods can reach a $\sim$4$\times$10$^{-3}$ floor;
  there GL needs roughly half the latitude points of DH (see
  \texttt{verify\_gl\_grid.py}).
\item The wall-clock benchmark above does NOT load pretrained weights; it
  times architecture-level forward cost. The cost depends only on
  architecture and tensor shapes, not on weight values.
\end{enumerate}

\section*{Reproduce}

\begin{lstlisting}
module load pytorch/2.6.0-1
# Equivariance for all configs at lmax=6, mmax=2
python3 experiments/verify_gl_grid.py
# Wall-clock: rigorous IID protocol, 5 runs x 30 distinct batches
PYTORCH_CUDA_ALLOC_CONF=expandable_segments:True \
  python3 experiments/expG_verify_speedup_v2.py \
    --batch_size 8 --n_runs 5 --n_iter 30 --n_warmup 10
# Drop-in test on actual trained checkpoints (small model)
PYTORCH_CUDA_ALLOC_CONF=expandable_segments:True \
  python3 experiments/expG_dropin_pretrained.py
\end{lstlisting}

\end{document}
